%%%%%%%%%%%%%%%%%%%%%%%%%%%%%%%%%%%%%%%%%%%%%%%%%%%%%%%%%%%%%%%%%%%%%%
%% Shared frame for the ESD build-up sequence in l02_esd.
%%
%% Four figures in the chapter are the same picture with more added each
%% time: a chip outline with three rails brought out to pads, then one
%% diode, then two, then the whole network with the grounded gate NMOS.
%% They only teach anything if they are literally the same drawing, so
%% the frame lives here and each figure adds its own devices.
%%
%% Coordinates: pads at x = 0, rails run to \esdRight. VDD at \esdVdd,
%% the signal pin at \esdPin, VSS at zero.
%%
%% The band between VSS and the pin has to hold a device, the gate
%% resistor beside it and a two line caption above it, which is what
%% sets \esdPin. The hand drawing these figures replace put every
%% device in that one band with the captions in a row above them, and
%% that is worth keeping: it is what lets the eye read the network as
%% one row of paths rather than four scattered symbols.
%%%%%%%%%%%%%%%%%%%%%%%%%%%%%%%%%%%%%%%%%%%%%%%%%%%%%%%%%%%%%%%%%%%%%%

\newcommand{\esdVdd}{5.8}
\newcommand{\esdPin}{3.4}
\newcommand{\esdRight}{9.8}

% Where a device in the lower band sits, and where its caption goes.
\newcommand{\esdLow}{1.7}
\newcommand{\esdCap}{2.45}

% A bond pad: the crossed square the chapter's originals use. The
% latch-up figure needs one away from the rails, so the general form
% takes a position and the rail form is written in terms of it.
\newcommand{\esdPadAt}[2]{%
  \draw ({#1-0.17},{#2-0.17}) rectangle ({#1+0.17},{#2+0.17});
  \draw ({#1-0.17},{#2-0.17}) -- ({#1+0.17},{#2+0.17});
  \draw ({#1-0.17},{#2+0.17}) -- ({#1+0.17},{#2-0.17});
}
\newcommand{\esdPad}[1]{\esdPadAt{0}{#1}}

% The chip outline, three rails, their pads and their names. #1 is the
% name of the middle rail, because the overview figure calls it PIN and
% the protection figures treat it as one signal pin among many.
\newcommand{\esdframe}{%
  \draw[dotted, gray] (-0.8,-1.3) rectangle ({\esdRight+0.3},{\esdVdd+1.3});
  \foreach \y/\name/\num in {\esdVdd/VDD/1, \esdPin/PIN/2, 0/VSS/0} {
    \draw (0,\y) -- (\esdRight,\y);
    \esdPad{\y}
    \node[armygreen, anchor=south west, scale=0.85] at (0.05,{\y+0.2}) {\name};
    %- above the rail, not left of it: the zap wiring comes in on the
    %- left and the numbers were landing on top of it
    \node[red, anchor=south east, scale=0.9] at (-0.25,{\y+0.18}) {\num};
  }
}

% A vertical diode conducting upward, from (#1,#2) to (#1,#3), with its
% body centred at #4 and captioned #5. Every protection device in these
% figures is one of these; drawing them through one macro is what makes
% the reverse-biased-in-normal-operation argument in the text hold
% across all four pictures.
%
% The centre is explicit because a diode spanning VSS to VDD would
% otherwise put its symbol exactly on the pin rail it is meant to cross
% without connecting.
%
% The caption goes above the body and to the left of the column, the
% way the hand drawing had it. Beside the body it collides with the
% next column along; above and centred it lands on the diode's own
% wire.
\newcommand{\esdDiode}[5]{%
  %- the house diode from ckt_lib, so ESD diodes match every other
  %- diode in the book
  \draw (#1,#2) -- (#1,{#4-0.8}) \vdiode{} -- (#1,#3);
  \fill (#1,#2) circle (0.075);
  \fill (#1,#3) circle (0.075);
  \node[red, anchor=south east, scale=0.8, align=right]
    at ({#1-0.15},{#4+0.85}) {#5};
}

% A grounded gate NMOS at x=#1, drain on the rail at #2, source on the
% rail at #3, body centred at #4, node id #5 (letters only - it names a
% TikZ node) and caption #6. The centre is explicit for the same reason
% the diode's is: left to the midpoint, a VDD to VSS device lands on the
% pin rail it is meant to cross without touching. The gate goes to the
% source through a resistor, which is the whole trick: the resistor does
% nothing at DC, so the device is off whenever the chip is running, but
% during a zap the substrate current it develops lifts the gate and
% turns the parasitic bipolar on.
%
% The resistor gets its own column a full unit to the left. Any closer
% and the zigzag is drawn on top of the transistor, which loses the one
% thing the figure is here to show. It is \vresistor from ckt_lib, so it
% matches every other resistor in the book; that fixes its height at
% 1.5, and the gate lead reaches down to meet it.
\newcommand{\esdggn}[6]{%
  \node[nmos] (esdm#5) at ({#1+0.35},#4) {};
  \draw (esdm#5.D) -- ({#1+0.35},#2);
  \draw (esdm#5.S) -- ({#1+0.35},#3);
  \fill ({#1+0.35},#2) circle (0.075);
  \fill ({#1+0.35},#3) circle (0.075);
  \draw (esdm#5.G) -- ({#1-1.0},#4);
  \draw ({#1-1.0},#4) -- ({#1-1.0},{#3+1.5});
  \draw ({#1-1.0},#3) \vresistor{};
  \fill ({#1-1.0},#3) circle (0.075);
  %- above the device and left of its column, level with the diode
  %- captions: beside the device it runs into the next column along,
  %- and centred above it lands on the drain wire
  \node[red, anchor=south east, scale=0.8, align=right]
    at ({#1+0.05},{#4+0.75}) {#6};
}

% The zap itself: current forced in at one pad, taken out at another.
\newcommand{\esdZapIn}[1]{%
  \draw (-2.7,#1) to [I] (-1.3,#1);
  \draw (-1.3,#1) -- (0,#1);
}
\newcommand{\esdZapOut}[1]{%
  \draw (0,#1) -- (-1.3,#1) -- (-1.3,{#1-0.5});
  \draw (-1.3,{#1-0.5}) \vground;
}

% ------------------------------------------------------------------
% Latch-up cross section for l02_esd figures 15/16: an inverter and a
% negative-going pin, drawn on the carrier colour language - red p
% regions, blue n regions, white depletion gaps, teal oxide. Figure 15
% shows the injected carriers; figure 16 adds the parasitic bipolars.
% Surface at y=0, x in [0,19].
% ------------------------------------------------------------------
\newcommand{\scrimplant}[5]{% x0 x1 depth colour label
  \fill[white, rounded corners=3pt] ({#1-0.18},0.02) rectangle ({#2+0.18},{#3-0.18});
  \fill[#4, rounded corners=2pt] (#1,0) rectangle (#2,#3);
  \draw[rounded corners=2pt] (#1,0) rectangle (#2,#3);
  \node[white, anchor=center, scale=0.85] at ({(#1+#2)/2},{#3/2}) {#5};
}

\newcommand{\scrframe}{%
  % p- substrate
  \fill[hcharge!14] (0,0) rectangle (19,-6.5);
  % n-well with a white depletion border
  \fill[white] (6.1,0.02) rectangle (13.9,-5.3);
  \fill[echarge!12] (6.4,0) rectangle (13.6,-5.0);
  \draw (6.4,0) -- (6.4,-5.0) -- (13.6,-5.0) -- (13.6,0);
  % surface
  \draw[line width=1.4pt] (0,0) -- (19,0);
  % implants: inverter NMOS on the left, PMOS in the well, well tap,
  % and the pin diffusion far right
  \scrimplant{0.5}{1.6}{-0.8}{hcharge!70}{$p+$}
  \scrimplant{2.4}{3.7}{-0.8}{echarge!75}{$n+$}
  \scrimplant{4.9}{6.0}{-0.8}{echarge!75}{$n+$}
  \scrimplant{7.4}{8.5}{-0.8}{hcharge!70}{$p+$}
  \scrimplant{9.7}{10.8}{-0.8}{hcharge!70}{$p+$}
  \scrimplant{11.6}{13.2}{-1.3}{echarge!75}{$n+$}
  \scrimplant{16.2}{17.8}{-0.8}{echarge!75}{$n+$}
  % gates: teal oxide with a white poly slab on top
  \fill[oxteal!60] (3.75,0) rectangle (4.85,0.16);
  \draw (3.75,0) rectangle (4.85,0.16);
  \draw[fill=white, rounded corners=3pt] (3.75,0.16) rectangle (4.85,0.75);
  \fill[oxteal!60] (8.55,0) rectangle (9.65,0.16);
  \draw (8.55,0) rectangle (9.65,0.16);
  \draw[fill=white, rounded corners=3pt] (8.55,0.16) rectangle (9.65,0.75);
  % metal: ground on the left pair, In on the gates, Out on the drains,
  % supply on the well tap, and the pin pulled below ground
  \draw (1.05,0) -- (1.05,0.5) -- (3.05,0.5) -- (3.05,0);
  \draw (0.45,0.5) -- (1.05,0.5);
  \draw (0.45,0.5) \vground;
  \draw (4.3,0.75) -- (4.3,2.1) -- (9.1,2.1) -- (9.1,0.75);
  \draw (6.0,2.1) -- (6.0,2.6);
  \fill (6.0,2.1) circle (0.075);
  \draw (6.0,2.69) circle (0.09);
  \node[anchor=south] at (6.0,2.85) {In};
  \draw (5.45,0) -- (5.45,1.1) -- (7.95,1.1) -- (7.95,0);
  \draw (7.3,1.1) -- (7.3,1.5);
  \fill (7.3,1.1) circle (0.075);
  \draw (7.3,1.59) circle (0.09);
  \node[anchor=west] at (7.45,1.62) {Out};
  \draw (10.25,0) -- (10.25,0.5) -- (12.4,0.5) -- (12.4,0);
  \draw (12.4,0.5) -- (12.4,1.1);
  \draw (12.4,1.1) \vsupply;
  \draw (17.0,0) -- (17.0,1.4);
  \draw (17.0,1.49) circle (0.09);
  \node[echarge, anchor=south, scale=1.0] at (17.0,1.65) {$V<0$};
  % region names
  \node[echarge!60!black, anchor=center] at (12.7,-4.5) {$n-$};
  \node[hcharge!60!black, anchor=center] at (16.2,-5.9) {$p-$};
}

% the injected carriers, shared by both figures
\newcommand{\scrarrows}{%
  \draw[echarge, line width=1.8pt, ->]
    (17.0,-0.9) .. controls (16.5,-2.8) and (15.2,-3.3) .. (13.85,-3.1);
  \node[draw, echarge, circle, scale=0.8, fill=white] at (15.9,-1.7) {1};
  \draw[hcharge, line width=1.8pt, ->]
    (10.2,-0.9) .. controls (9.7,-3.6) and (6.0,-4.6) .. (3.2,-4.2);
  \node[draw, hcharge, circle, scale=0.8, fill=white] at (8.1,-3.4) {2};
  \draw[echarge, line width=1.8pt, ->]
    (3.1,-0.9) .. controls (3.6,-2.4) and (5.2,-2.8) .. (6.6,-2.55);
  \node[draw, echarge, circle, scale=0.8, fill=white] at (4.5,-1.6) {3};
}
